\subsection{CW-Complexes}

\begin{proposition}
    Let $(X,A)$ be a relative CW-complex. If $A$ is Hausdorff (resp. normal) then $X$ is Hausdorff (resp. normal) as well. In particular every absolute CW-complex is Hausdorff and normal.
\end{proposition}
\begin{proof}
    We show the statement for the case that $A$ is normal. The Hausdorff case follows completely analogously by replacing the sets to be seperated by point sets. So let $B,C \subseteq X$ be closed disjoint sets. We will show the following claim inductively:
    \begin{claim}
    There exist open disjoint sets $U_n,V_n \subseteq X_n$, such that $B_n := B \cap X_n \subseteq U_n$ and $C_n := C \cap X_n \subseteq V_n$ as well as $U_{n-1} \subseteq U_n, V_{n-1} \subseteq V_n$.
    \end{claim}
    The case $n=-1$ follows from the assumption on $A$. For the induction step choose a push-out
    \begin{center}
        \commsquare{$\coprod_{i \in I} S^{n-1}$}{$X_{n-1}$}{$\coprod_{i \in I} D^n$}{$X_n$}{$(q_i)$}{}{}{$(Q_i)$}
    \end{center}
    and set $U_{n-1}^\prime := (q_i)^{-1}_{i \in I}(U_{n-1})$, $V_{n-1}^\prime := (q_i)^{-1}_{i \in I}(V_{n-1})$. By Lemma \ref{lem:technical_lemma1} shown below we can find open neighbourhoods $U_n^\prime,V_n^\prime \subseteq \coprod_I D^n$ with the following properties:
    \begin{enumerate}[(i)]
        \item $(Q_i)_I^{-1}(B_n) \subseteq U_n^\prime$ and $(Q_i)_I^{-1}(C_n) \subseteq V_n^\prime$
        \item $U_n^\prime \cap V_n^\prime = \emptyset$
        \item $U_{n-1}^\prime = U_n^\prime \cap \coprod_I S^{n-1}$ and $V_{n-1}^\prime = V_n^\prime \cap \coprod_I S^{n-1}$ 
    \end{enumerate}
    Then define 
    \[
    U_n := U_{n-1} \cup (Q_i)_I(U_n^\prime), \quad V_n := V_{n-1} \cup (Q_i)_I(V_n^\prime)
    \]
    By (iii) $U_n \cap X_{n-1} = U_{n-1}$ and $(Q_i)_I^{-1}(U_n) = U_n^\prime$, so $U_n \subseteq X_n$ is open. The same holds for $V_n$. Furthermore, $B_n \subseteq U_n, \; C_n \subseteq V_n$ by (i) and $U_n \cap V_n = \emptyset$ by (ii). This shows the claim. Now consider
    \[
    U := \bigcup_n U_n, \quad V := \bigcup_n V_n.
    \] 
    By construction $U \cap X_n \subseteq$ is open for all $n \geq -1$ and hence $U \subseteq X$ is open by definition of the colimit topology. The same holds for $V$. Moreover, $B \subseteq U, C \subseteq V$ and $U \cap V = \emptyset$. 
\end{proof}

It remains to proof the lemma mentioned in the proof of the above proposition.

\begin{lemma}
    \label{lem:technical_lemma1}
    Let $B , C \subseteq D^n$ be closed disjoint sets and $U^\prime, V^\prime \subseteq S^{n-1}$ open disjoint neighbourhoods of $B \cap S^{n-1}, C \cap S^{n-1}$, respectively. Then there exist open sets $U,V \subseteq D^n$ with the following properties:
    \begin{enumerate}[(i)]
        \item $B \subseteq U, C \subseteq V$
        \item $U \cap V =\emptyset$
        \item $U^\prime = U \cap S^{n-1}, V^\prime = V \cap S^{n-1}$
    \end{enumerate}
\end{lemma}

\begin{proof}
    Consider the open disjoint sets
    \[
        U^{\prime\prime} := \left\{ rx \mid x \in U^\prime,\, 0<r\leq 1\right\}
    \]
    and
    \[
        V^{\prime\prime} := \left\{ rx \mid x \in V^\prime,\, 0<r\leq 1\right\}.
    \]
    Moreover, since $B$ and $\overline{V}^\prime$ are disjoint and closed we can find an open neighbourhood $W_B$ of $B$, such that 
    \[
    \overline{W}_B \cap \overline{V}^\prime = \emptyset.
    \]
    Similialry we can find an open neighbourhood $W_C$ of $C$ with 
    \[
        \overline{W}_C \cap \overline{U}^\prime = \emptyset.
    \]
    After potentially shrinking $W_B$ and $W_C$ we may also assume, that they are disjoint. Now set 
    \[
        U := (U^{\prime\prime} - \overline{W}_C) \cup (W_B - S^{n-1})
    \]
    and 
    \[
        V := (V^{\prime\prime} - \overline{W}_B) \cup (W_C - S^{n-1}).
    \]
    It is then easy to verify, that these sets satisfy the statement.
\end{proof}